\DocumentMetadata{
  pdfversion=2.0,
  pdfstandard=ua-2,
  lang=en-US,
  tagging=on,
  tagging-setup={math/setup=mathml-SE,tabsorder=structure}
}
\documentclass[11pt,letterpaper]{article}
\usepackage[margin=0.68in,includefoot]{geometry}
\usepackage{newcomputermodern}
\usepackage{amsmath,amscd,mathtools}
\usepackage{tikz,adjustbox}
\providecommand{\square}{\mdlgwhtsquare}
\usepackage[hidelinks]{hyperref}
\hypersetup{
  pdftitle={Combinatorics PhD Exam, January 2012},
  pdfauthor={University of Florida Department of Mathematics},
  pdfsubject={Graduate mathematics examination},
  pdfdisplaydoctitle=true,
  bookmarksopen=true
}
\setcounter{secnumdepth}{0}
\setlength{\parindent}{0pt}
\setlength{\parskip}{0.28em}
\setlength{\emergencystretch}{3em}
\setlength{\leftmargini}{2.15em}
\setlength{\leftmarginii}{2.65em}
\setlength{\leftmarginiii}{3em}
\renewcommand{\labelenumi}{\textbf{\theenumi.}}
\pagestyle{empty}
\raggedbottom
\allowdisplaybreaks
\newenvironment{examproblems}[1]{%
  \begin{enumerate}%
  \setcounter{enumi}{#1}%
  \setlength{\topsep}{0.35em}%
  \setlength{\partopsep}{0pt}%
  \setlength{\itemsep}{0.48em}%
  \setlength{\parsep}{0.18em}%
}{\end{enumerate}}
\newenvironment{examsubparts}{%
  \begin{enumerate}%
  \setlength{\topsep}{0.18em}%
  \setlength{\partopsep}{0pt}%
  \setlength{\itemsep}{0.16em}%
  \setlength{\parsep}{0.1em}%
}{\end{enumerate}}
\newenvironment{examinstructions}{%
  \par\begingroup\itshape
}{%
  \par\endgroup\vspace{0.18em}
}
\makeatletter
\renewcommand\section{\@startsection{section}{1}{0pt}%
  {-1.25ex plus -.25ex minus -.1ex}%
  {0.55ex plus .1ex}%
  {\normalfont\large\bfseries}}
\renewcommand{\@maketitle}{%
  \newpage\null\vskip -1em%
  {\footnotesize\bfseries
    \href{https://www.ufl.edu/}{University of Florida}, \href{https://math.ufl.edu/}{Department of Mathematics}\par}%
  \vskip -0.5em%
  \tagstructbegin{tag=Title}%
    {\LARGE\bfseries\href{https://gma.math.ufl.edu/exams/combinatorics-phd/}{Combinatorics PhD Exam}, January 2012\par}%
  \tagstructend%
  \vskip 0.25em%
  \tagmcbegin{artifact}%
    \hrule height 1pt%
  \tagmcend%
  \vskip 1em%
}
\makeatother
\title{Combinatorics PhD Exam, January 2012}
\author{}
\date{}
\begin{document}
\maketitle
\thispagestyle{empty}
\begin{examproblems}{0}
\item
Prove that there exists a simple~\(k\)-regular graph on~\(n\) vertices if and only if \(nk\) is even, and \(n \geq k+1\).
\item
Find a closed formula (no summation signs) for the generating function \(\sum_{n\geq 0} g(n)x^n/n!\), where \(g(n)\) is the number of permutations of length~\(n\) with all cycle lengths even.
\item
Show that a planar graph in which every face has an even number of vertices must be bipartite.
\item
Let \(f(n)\) be the number of simple graphs on vertex set \(\{1,2,\ldots,n\}\) in which each connected component is a cycle. Find a closed form for the generating function \(\sum_{n\geq 0} f(n)x^n/n!\).
\item
How many Hamiltonian cycles does the graph \(K_{n,n}\) have?
\item
Let \(h(n)\) be the number of compositions of the integer~\(n\) into odd parts. Find the generating function \(H(x)=\sum_{n\geq 1}h(n)x^n\).
\item
Find a formula for the number of partitions of the set \(\{1,2,\ldots,n\}\) in which each block contains more than one element. Your answer can contain the symbol \(B(i)\) for the number of all partitions of \([i]\).
\item
Let \(d\geq 2\) be an integer, and let~\(P\) be the poset of all vectors of length~\(d\) with non-negative integer coordinates in which \(\mathbf{x}\leq\mathbf{y}\) if \(x_i\leq y_i\) for all~\(i\).
\begin{examsubparts}
\item[\textnormal{(a)}]
Let~\(n\) be any positive integer. Does~\(P\) contain an antichain of~\(n\) elements?
\item[\textnormal{(b)}]
Does~\(P\) contain an infinite antichain?
\end{examsubparts}
\end{examproblems}
\end{document}
